\begin{lemma}[Defective-chain structure, part (e)]
Let $(T,e_0,\varphi)$ be a globally maximum triple for $G$.  Let
$\alpha\in\mathcal D(V(T),\varphi)$ be defective, let
$\gamma\in\mathcal F(T,e_0,\varphi)$ be free and missing at a vertex
$u\in V(T)$, and let $z\notin V(T)$.  Let the supplied incidence walk $P$
from $u$ to $z$ be a full defective-chain path for
$(T,\varphi,\alpha,\gamma,u,z)$.  If the supplied indices $i_1<i_2$ and
vertices $v_1,v_2$ satisfy the first-departure/last-return predicate for
$(T,\varphi,\alpha,\gamma,u,z,P)$, then $V(T)\cup\{v_1\}$ is elementary
with respect to $\varphi$.
\end{lemma}
\begin{proof}
Put $X=V(T)$.  Suppose that a color $\delta$ is missing both at $v_1$ and
at some $x\in X$.  By
\noderef{ScheideDefectiveChainBoundaryEdges} and the three-boundary-edge
clause of \noderef{ScheideGlobalMaximumBoundaryFacts}, the path has at
least two boundary crossings after its first departure.  Thus $v_1$ is
not an end of the full $(\alpha,\gamma)$-path, so both $\alpha$ and
$\gamma$ are present at $v_1$ and
$\delta\notin\{\alpha,\gamma\}$.

Consider the full $(\gamma,\delta)$-component $C$ through $v_1$.  Both
colors are missing on the closed set $X$; hence no edge identity of either
color leaves $X$, by \noderef{ScheideGlobalMaximumBoundaryFacts}.  The
outside component $C$ is therefore disjoint from $X$.  Switch $C$.
Properness of the original coloring follows from the $k$-triple clause of
\noderef{GlobalMaximumTriple}, so it is preserved by
\noderef{TwoColorComponentSwapProper}; here $\gamma\ne\delta$ was proved
in the preceding paragraph.  No tree edge or tree missing-color witness
changes, and $\gamma$ remains missing at $u$.  At $v_1$, $\gamma$ is now
missing, and the unique
$\gamma$-edge of the old path following $v_1$ has become a
$\delta$-edge.  Hence the initial segment $uPv_1$ is now the full
$(\alpha,\gamma)$-component through $u$.

Switch that initial component.  The colors $\alpha,\gamma$ are unused on
$T$ by \noderef{ScheideGlobalMaximumUnused}, so the same sequence remains
Tashkinov, and $\alpha$ becomes missing at $u$.  The colors are distinct,
since $\gamma$ is missing at $u$ whereas the defective color $\alpha$ is
present there by the two corresponding clauses of
\noderef{ScheideGlobalMaximumBoundaryFacts}; hence properness of this
second switch again follows from \noderef{TwoColorComponentSwapProper}.
This initial segment uses only the first $\alpha$-boundary edge of $P$.
At least two later $\alpha$-boundary edge identities retain their colors, by
\noderef{ScheideDefectiveChainBoundaryEdges}.  Appending either to the
tree contradicts global maximality.  Thus a color missing at $v_1$ cannot
be missing on $X$.  Since $X$ itself is elementary by
\noderef{ScheideGlobalMaximumOrder}, the union $X\cup\{v_1\}$ is
elementary.
\end{proof}
